\section{GPT-6 Astraによる物理世界の設計支援}
OpenAIが、ChatGPTとCodex、特にGPT-6 Astraを使った機械設計・電気設計・シミュレーションの事例を紹介した。公式投稿によれば、ハードウェアチームのトークン使用量は2月比23倍となり、MazdaやChip Ganassi Racingとの協業例も挙げられた。2人のチームメンバーは、Codexと対話するプロトタイプ装置を1週間未満で設計したという。機能要件、筐体、PCB、ソフトウェア統合をAIが支援した一方、製造には複数回のレビューを要した。BenchCAD 2.0の評価値も示されたが、外部での再現や製造業全般への適用可能性は未確認である。
\dailyxsource{OpenAI、11月4日03:56 JST}{https://x.com/OpenAI/status/2117692475775258785}

\section{Anthropicが推論証明技術FlashProveを公開}
Anthropicは、応答が指定されたモデルから生成されたことを暗号学的に証明する「FlashProve」をオープンソース公開したと発表した。Grokが収集した研究紹介では、zkSNARKを利用し、4-bitのLlama 4 Maverick（400B）について単一H100で約110トークン毎秒の証明生成を報告。証明サイズ約19MB、CPU検証約17秒とされる。ただし通常推論に比べGPU時間は70～95倍必要とされ、運用負荷は小さくない。コードとLean証明の公開も紹介されたが、第三者による独立再現は未確認である。
\dailyxsource{Anthropic、11月4日05:10 JST}{https://x.com/AnthropicAI/status/2117710942826123294}

\section{DeepMind英国従業員の契約条項を巡る報道}
Washington Postの記事を引用する投稿が、Google DeepMindの英国従業員による契約条項への反発を伝えた。報道によれば、退職後も公的な批判を生涯禁じるとされる条項が、AIリスクに関する警告を妨げかねないとの懸念がある。X上では社内議論をめぐる表現も拡散したが、Grokが確認したのは記事を引用した二次投稿であり、従業員本人の一次投稿や契約条項の原文、Google側の正式な応答は未確認である。報道内容と確認済み事実を区別して読む必要がある。
\dailyxsource{Choblin、11月4日06:58 JST}{https://x.com/choblin29/status/2117738331697983813}

\section{AMD決算とHeliosを巡るAI需要}
AMDの第3四半期決算を引用する投稿では、データセンター事業の売上高が約85億ドルとなり、前年同期比でほぼ倍増したと紹介された。また、Anthropic、Meta、Microsoft、OpenAI、OracleなどがHeliosラックスケールシステム上でモデルを稼働させているとの言及も共有された。AI計算需要の強さを示す材料として注目される一方、Grokの収集では決算リリースとの独立照合は行われていない。数値の定義や具体的なワークロードは、公式資料での確認が残る。
\dailyxsource{PreReason、11月4日06:58 JST}{https://x.com/Pre_Reason/status/2117738312936911180}

\section{次期GPT・Claudeなどの未確認情報}
11月にGPT-6.2 Astra/Bel、Claude Mythos/Fable 5.5、Grokの更新などが登場するという噂が、複数のXアカウントで共有された。API上でモデル名が一時的に見えたとの報告もあるが、Grokの観測では公式発表や一次資料を確認できていない。似た一覧が複数投稿されたことは製品の存在や公開時期を裏付けるものではない。現時点では未確認のリーク・予測として扱い、正式なリリース情報とは明確に区別する。
\dailyxsource{anya、11月3日12:50 JST}{https://x.com/anya_yfa7/status/2117464475775516956}

\section{法律特化GLMモデルの提供を巡る投稿}
Zhipuの「GLM-5.3-FlashX-Legal」がOpenRouterで利用できるようになったとする投稿があった。法律タスク向けの追加学習を施したモデルとされ、100万トークンのコンテキスト長や価格、Harvey LABでの高いall-pass率などが紹介された。エージェントとMCPを組み合わせた法律実務への利用も提案されている。ただし、GrokはOpenRouter上のモデルの存在やベンチマーク結果を独立確認していない。提供条件と各法域での性能評価は今後の確認事項となる。
\dailyxsource{Jorge Araujo、11月4日06:57 JST}{https://x.com/JorgeAraujo/status/2117737878637642210}

\section{米金融機関がAIエージェントのログインを制限か}
米国の銀行・証券会社が第三者AIエージェントによる認証情報の利用を制限しているとの投稿が共有された。法律事務所のアップデートを引用した内容では、Chase、Charles Schwab、Citiなどがエージェントによるログインを遮断し、Wells Fargoは利用者責任に関する条項を加えたとされる。自律エージェントと金融口座の接続に伴うセキュリティ・責任分担が論点だが、各社の公式声明や規約原文はGrokの観測では未確認である。
\dailyxsource{Auditrol、11月4日06:54 JST}{https://x.com/auditrol/status/2117737258287591856}

\section{Higgsfieldなど生成AIツールの更新情報}
生成AIツールの新機能をまとめた投稿では、Higgsfieldがアイデアからマンガ全体を生成できるようになったと紹介された。動画付きのスレッドで複数のツール更新を取り上げており、画像・映像制作の工程を一括で支援する方向性への関心がうかがえる。ただし、Grokが確認したのは紹介投稿であり、各ツールの公式リリース情報や機能範囲は未確認である。実際の提供条件や出力品質については追加確認が必要となる。
\dailyxsource{Dheepan Ratnam、11月3日20:30 JST}{https://x.com/Dheepanratnam/status/2117580224527491408}
